% !TEX root = ../../LAFROX_AnalisisFinanciero_2.tex
\section{Valorización de las cinco innovaciones}
\label{sec:e2-innovaciones}

Esta sección valoriza las cinco innovaciones de la Oferta Técnica según el apartado 2.7 del Formulario E-21: inversión, costo operacional incremental y beneficio esperado de cada una. Los resultados económicos de esta sección corresponden a la hoja \texttt{INNOVACIONES AGRUPADAS} del Modelo Financiero (Entregable 3), que agrupa las horas y los costos por innovación y por clase (inversión u operación). El Formulario T-15 respalda las partidas y el esfuerzo (LafroX, 2026c, sección 6.3). El Formulario T-19 define los indicadores de beneficio (LafroX, 2026d, fichas 1 a 5) y el Anexo 4-W aporta el dimensionamiento de nube (LafroX, 2026a, tablas 4-W.I1 y 4-W.I2).

Los montos se expresan en UF. Los totales por innovación y por período se presentan, además, en CLP y USD, con los valores de la hoja \texttt{PARAMETROS}: UF de \$40.000 y dólar de \$900 (Distribuidora Puelche S.A., 2026a, Formulario E-24).

\begin{tablalafrox}{Innovaciones valorizadas}{tab:e2-inn-lista}{L{1.6cm} Y}
  {\cab{Código} & \cab{Innovación}}
  INN-01 & Seguimiento de vencimiento posentrega en el local \\
  INN-02 & Reproducción de incidentes de terreno \\
  INN-03 & Vida útil remanente por historia térmica del lote \\
  INN-04 & Tramo variable de la Operación ligado al costo de servir \\
  INN-05 & Hoja de negocio del almacenero \\
\end{tablalafrox}

% ---------------------------------------------------------------------------
\subsection{Criterio de valorización}
\label{sec:e2-inn-criterio}

La valorización contiene tres conceptos con su tratamiento correspondiente:

\begin{itemize}
  \item \textbf{Costo de LafroX.} Corresponde a las HH del equipo y a los servicios de nube que cada innovación requiere. Se divide en inversión, ejecutada durante la implementación, y costo operacional incremental, durante la Operación.
  \item \textbf{Beneficio para el cliente.} Es el beneficio estimado para Puelche o el almacenero según el indicador. Se presenta por separado del flujo de LafroX.
  \item \textbf{Ingreso variable de LafroX.} Solo INN-04 genera un ingreso para el proponente: el tramo variable de la Operación ligado al costo de servir.
\end{itemize}


% ---------------------------------------------------------------------------
\subsection{Datos de entrada y supuestos}
\label{sec:e2-inn-datos}

La sección A de la hoja proporciona los datos que alimentan el cálculo de la hoja de innovaciones.

\begin{tablalafrox}{Datos de entrada del cálculo de innovaciones}{tab:e2-inn-datos}{L{1.2cm} Y R{2.6cm} L{2.6cm} L{2.2cm}}
  {\cab{Ref.} & \cab{Dato} & \cab{Valor} & \cab{Unidad} & \cab{Tipo}}
  D-01 & Unidades despachadas al mes & 2.400.000 & unidades/mes & Documento \\
  D-02 & Pedidos al mes & 31.000 & pedidos/mes & Documento \\
  D-03 & Visitas de preventa al mes & 62.000 & visitas/mes & Documento \\
  D-04 & Recepciones de proveedor al mes & 1.150 & recepciones/mes & Documento \\
  D-05 & Ventas anuales & 148.000.000.000 & CLP/año & Documento \\
  D-06 & Clientes del canal tradicional & 11.600 & locales & Documento \\
  D-07 & Clientes activos & 14.200 & clientes & Documento \\
  D-08 & SKU refrigerados o congelados & 1.100 & SKU & Documento \\
  D-09 & SKU activos & 8.400 & SKU & Documento \\
  D-10 & INN-02: incidentes por semana & 1 & incidentes/semana & Supuesto \\
  D-11 & INN-02: personas detenidas por incidente & 2 & personas & Supuesto \\
  D-12 & Valor del minuto-persona de reparto & 0,002 & UF/min & Supuesto \\
  D-13 & Valor del minuto del almacenero & 0,001 & UF/min & Supuesto \\
  D-14 & INN-04: tramo variable máximo mensual & 10 & UF/mes & Supuesto \\
  D-15 & INN-04: primer mes con pago del tramo & 24 & mes & Documento \\
\end{tablalafrox}

% ---------------------------------------------------------------------------
\subsection{Inversión y costo operacional en horas hombre}
\label{sec:e2-inn-hh}

La sección B de la hoja asigna las HH de cada innovación y clase a un perfil que las valoriza, con su costo por HH. Cada fila tiene un solo perfil y un tramo de meses, con las HH repartidas en partes iguales entre esos meses.

\begin{tablalafrox}{HH de innovación por innovación y clase}{tab:e2-inn-hh}{L{1.4cm} L{1.7cm} R{1.1cm} R{1.3cm} L{1.5cm} Y R{1.5cm}}
  {\cab{Innov.} & \cab{Clase} & \cab{Filas} & \cab{HH} & \cab{Meses} & \cab{Perfil que valoriza} & \cab{UF}}
  INN-01 & Inversión & 1 & 960,0 & 7–12 & Ingeniero de Datos y Bases Senior & 768,0 \\
  INN-02 & Inversión & 1 & 960,0 & 9–10 & Analista QA Senior & 480,0 \\
  INN-02 & Operación & 22 & 576,0 & 22–55 & Analista QA Senior, Analista QA Junior y Líder de Calidad y Pruebas & 272,2 \\
  INN-03 & Inversión & 1 & 960,0 & 7–12 & Ingeniero de Datos y Bases Senior & 768,0 \\
  INN-03 & Operación & 11 & 288,0 & 24–54 & Ingeniero de Datos y Bases Senior y Junior, y Líder de Datos & 230,4 \\
  INN-04 & Inversión & 1 & 160,0 & 17–23 & Jefe de Proyecto & 232,0 \\
  INN-04 & Operación & 1 & 528,0 & 24–56 & Jefe de Proyecto & 765,6 \\
  INN-05 & Inversión & 2 & 1.120,0 & 13–23 & Coordinador de Implantación & 672,0 \\
  \midrule
  \textbf{Total} & & \textbf{40} & \textbf{5.552,0} & & & \textbf{4.188,2} \\
\end{tablalafrox}

Las 5.552 HH equivalen al 2,5\,\% de las HH del proyecto. La operación de INN-02 e INN-03 ocupa varias filas porque los perfiles de calidad y de datos solo tienen HH en algunos meses de la Operación: cada fila corresponde a un mes con dotación disponible.

% ---------------------------------------------------------------------------
\subsection{Costos sin horas hombre: nube y otros}
\label{sec:e2-inn-nube}

La sección C de la hoja estima el consumo de servicios de nube de cada innovación, a partir del dimensionamiento de nube del SD4 para la región de São Paulo.

\begin{tablalafrox}{Costos sin HH de las innovaciones}{tab:e2-inn-nube}{L{1.1cm} Y L{1.3cm} L{1.5cm} R{1.3cm} R{1.4cm} L{2.3cm}}
  {\cab{Ref.} & \cab{Rubro} & \cab{Innov.} & \cab{Meses} & \cab{UF/mes} & \cab{Total UF} & \cab{Estado}}
  C-01 & Cálculo nocturno y almacenamiento N-10 & INN-01 & 16–56 & 0,23 & 9,57 & Estimación base \\
  C-02 & Cómputo efímero para reproducción en QA y Preproducción & INN-02 & 5–56 & 0,01 & 0,43 & Estimación base \\
  C-03 & Cómputo y resultados en N-10 & INN-03 & 16–56 & 0,81 & 33,10 & Estimación base \\
  C-04 & Hardware de borde & INN-03 & 16–56 & 0,00 & 0,00 & Hardware ofertado \\
  C-05 & Laboratorio externo para validar lotes & INN-03 & 11–15 & 0,00 & 0,00 & Sin contratación \\
  C-06 & Cálculo y almacenamiento N-10 (ciclo mensual) & INN-05 & 21–56 & 0,03 & 1,16 & Estimación base \\
  \midrule
  \textbf{Total} & & & & & \textbf{44,26} & \\
\end{tablalafrox}

El hardware de borde de INN-03 ya figura en el Formulario T-11, por lo que no tiene compra adicional. El laboratorio externo solo se contrataría si el área de Calidad lo pide, y en el escenario base no se contrata.
% ---------------------------------------------------------------------------
\subsection{Beneficio esperado para el cliente}
\label{sec:e2-inn-beneficio}

La sección D de la hoja calcula el beneficio mensual de cada innovación como volumen por mes $\times$ tasa base $\times$ reducción $\times$ valor unitario.
\begin{tablalafrox}{Beneficio esperado para el cliente}{tab:e2-inn-beneficio}{L{1.0cm} Y R{1.6cm} R{1.2cm} R{1.1cm} R{1.4cm} R{1.3cm} R{1.6cm}}
  {\cab{Ref.} & \cab{Indicador} & \cab{Volumen/mes} & \cab{Tasa base} & \cab{Reduc.} & \cab{Valor unit. UF} & \cab{Meses} & \cab{Total UF}}
  B-01 & INN-01: unidades vencidas evitadas & 2.400.000 & 1\,\% & 25\,\% & 0,1285 & 16–56 & 31.604,2 \\
  B-02 & INN-02: minutos de interrupción evitados & 4,3 & 120 min & 30\,\% & 0,004 & 13–56 & 27,5 \\
  B-03 & INN-03: lotes descartados evitados & 150,6 & 5\,\% & 20\,\% & 268,1159 & 16–56 & 16.554,6 \\
  B-04 & INN-04: minutos-persona de servir evitados & 31.000 & 20 min & 5\,\% & 0,002 & 24–56 & 2.046,0 \\
  B-05 & INN-05: minutos de reposición ahorrados & 50.647,9 & 10 min & 20\,\% & 0,001 & 21–56 & 3.646,6 \\
  \midrule
  \textbf{Total} & & & & & & & \textbf{53.878,8} \\
\end{tablalafrox}

Cada indicador se construye así:

\begin{itemize}
  \item \textbf{B-01.} El volumen corresponde a las unidades despachadas (D-01). El valor unitario es el precio de venta medio por unidad: ventas anuales divididas por las unidades del año (D-05 / (D-01 $\times$ 12)), expresado en UF. Es una cota superior y bruta, porque los canjes no están cuantificados.
  \item \textbf{B-02.} El volumen son los incidentes por mes (D-10 llevado a base mensual). El valor es el minuto-persona de reparto por las personas detenidas en cada incidente (D-12 $\times$ D-11).
  \item \textbf{B-03.} El volumen son las recepciones con SKU refrigerados o congelados (D-04 $\times$ D-08 / D-09), con la convención de un lote por recepción, para estimar el volumen de lotes. El valor es la venta media por recepción (D-05 / (D-04 $\times$ 12)), en UF. También es una cota superior.
  \item \textbf{B-04.} El volumen son todas las entregas (D-02), como cota superior. El valor es el minuto-persona de reparto (D-12).
  \item \textbf{B-05.} El volumen son las visitas de preventa al canal tradicional (D-03 $\times$ D-06 / D-07). El valor es el minuto del almacenero (D-13).
\end{itemize}


% ---------------------------------------------------------------------------
\subsection{Ingreso variable de LafroX (INN-04)}
\label{sec:e2-inn-ingreso}

La sección E de la hoja calcula el tramo variable máximo de la Operación asociado a INN-04: 10 UF por mes (D-14), desde el mes 24 (D-15) hasta el mes 56, es decir, 33 meses y 330 UF en total.

% ---------------------------------------------------------------------------
\subsection{Resumen por innovación}
\label{sec:e2-inn-resumen}

La sección F de la hoja consolida, para cada innovación, la inversión, el costo operacional incremental y el beneficio esperado que exige el apartado 2.7 del E-21.

\begin{tablalafrox}{Resumen de la valorización por innovación (UF)}{tab:e2-inn-resumen}{L{1.3cm} R{1.6cm} R{1.6cm} R{1.5cm} R{1.6cm} >{\RaggedLeft\arraybackslash}X R{1.4cm}}
  {\cab{Innov.} & \cab{Inversión} & \cab{Costo op. HH} & \cab{Costo op. nube} & \cab{Egreso total} & \cab{Beneficio cliente} & \cab{Ingreso LafroX}}
  INN-01 & 768,0 & 0,0 & 9,6 & 777,6 & 31.604,2 & 0,0 \\
  INN-02 & 480,0 & 272,2 & 0,4 & 752,6 & 27,5 & 0,0 \\
  INN-03 & 768,0 & 230,4 & 33,1 & 1.031,5 & 16.554,6 & 0,0 \\
  INN-04 & 232,0 & 765,6 & 0,0 & 997,6 & 2.046,0 & 330,0 \\
  INN-05 & 672,0 & 0,0 & 1,2 & 673,2 & 3.646,6 & 0,0 \\
  \midrule
  \textbf{Total} & \textbf{2.920,0} & \textbf{1.268,2} & \textbf{44,3} & \textbf{4.232,4} & \textbf{53.878,8} & \textbf{330,0} \\
\end{tablalafrox}
\newpage
\begin{tablalafrox}{Egreso total y beneficio por innovación en tres monedas}{tab:e2-inn-monedas}{L{1.3cm} R{1.5cm} R{1.9cm} R{1.8cm} R{1.7cm} >{\RaggedLeft\arraybackslash}X R{1.9cm}}
  {\cab{Innov.} & \cab{Egreso UF} & \cab{Egreso MM\$} & \cab{Egreso USD} & \cab{Benef. UF} & \cab{Benef. MM\$} & \cab{Benef. USD}}
  INN-01 & 777,6 & 31,1 & 34.559 & 31.604,2 & 1.264,2 & 1.404.630 \\
  INN-02 & 752,6 & 30,1 & 33.449 & 27,5 & 1,1 & 1.220 \\
  INN-03 & 1.031,5 & 41,3 & 45.844 & 16.554,6 & 662,2 & 735.758 \\
  INN-04 & 997,6 & 39,9 & 44.338 & 2.046,0 & 81,8 & 90.933 \\
  INN-05 & 673,2 & 26,9 & 29.918 & 3.646,6 & 145,9 & 162.073 \\
  \midrule
  \textbf{Total} & \textbf{4.232,4} & \textbf{169,3} & \textbf{188.108} & \textbf{53.878,8} & \textbf{2.155,2} & \textbf{2.394.615} \\
\end{tablalafrox}

En total, las cinco innovaciones requieren 2.920,0 UF de inversión (\$116.800.000; USD 129.778) y 1.312,4 UF de costo operacional incremental (\$52.497.249; USD 58.330). El beneficio esperado para el cliente suma 53.878,8 UF (\$2.155.153.362; USD 2.394.615). En el escenario de pago máximo, el ingreso variable de INN-04 alcanza 330,0 UF (\$13.200.000; USD 14.667).

% ---------------------------------------------------------------------------
\subsection{Egreso e ingreso mensual de las innovaciones}
\label{sec:e2-inn-mensual}

La sección G de la hoja reparte los egresos y el ingreso de las innovaciones en los 56 meses del contrato. Para cada mes, suma las filas de las secciones B y C cuyo tramo incluye ese mes. Sus controles verifican que la suma mensual coincide con el resumen por innovación, tanto en el egreso como en el ingreso.

\begin{tablalafrox}{Egreso e ingreso de las innovaciones por período (UF)}{tab:e2-inn-periodos}{Y R{2.1cm} R{2.1cm} R{2.1cm} R{2.1cm}}
  {\cab{Concepto} & \cab{Meses 1–16} & \cab{Meses 17–20} & \cab{Meses 21–56} & \cab{Total}}
  Egreso HH de innovación & 2.444,8 & 307,8 & 1.435,6 & 4.188,2 \\
  Egreso de nube y otros & 1,1 & 4,2 & 38,9 & 44,3 \\
  \textbf{Egreso total de innovaciones} & \textbf{2.446,0} & \textbf{312,0} & \textbf{1.474,5} & \textbf{4.232,4} \\
  Ingreso variable INN-04 & 0,0 & 0,0 & 330,0 & 330,0 \\
\end{tablalafrox}
\newpage
\begin{tablalafrox}{Egreso total de las innovaciones por período en tres monedas}{tab:e2-inn-periodos-monedas}{Y R{2.3cm} R{2.6cm} R{2.1cm}}
  {\cab{Período} & \cab{UF} & \cab{CLP} & \cab{USD}}
  Meses 1–16 & 2.446,0 & 97.838.882 & 108.710 \\
  Meses 17–20 & 312,0 & 12.478.218 & 13.865 \\
  Meses 21–56 & 1.474,5 & 58.980.148 & 65.533 \\
  \midrule
  \textbf{Total} & \textbf{4.232,4} & \textbf{169.297.249} & \textbf{188.108} \\
\end{tablalafrox}

El 58\,\% del egreso de las innovaciones se concentra en los meses 1 a 16, donde ocurre la mayor parte de la inversión. En los meses 21 a 56, el egreso corresponde, sobre todo, al costo operacional en HH de INN-04 (765,6 UF), INN-02 (272,2 UF) e INN-03 (230,4 UF).

% ---------------------------------------------------------------------------
\subsection{Supuestos y limitaciones}
\label{sec:e2-inn-supuestos}

La lectura de los resultados depende de los supuestos de asignación y del escenario de beneficio.

\begin{itemize}
  \item \textbf{Perfil que valoriza.} El T-15 asigna HH por partida y perfil. La hoja económica utiliza su propia agrupación por perfil y mes. Para la conciliación se contrastan las filas de la hoja con la asignación del T-15 y sus ventanas.
  \item \textbf{Ubicación en el tiempo.} Los meses de las tablas corresponden a la imputación de la hoja económica. Se contrastan con las ventanas técnicas del T-15 y el T-19 antes de trasladar los montos al flujo consolidado.
  \item \textbf{Beneficios como cota superior.} B-01, B-03 y B-04 usan el volumen o el valor máximo usado en el escenario, por lo que representan el beneficio máximo esperado.
  \item \textbf{Nube por conciliar.} Los costos de la sección C son una estimación bruta con tarifas públicas, pendiente de conciliar con el consumo ya cubierto en Adquisiciones.
  \item \textbf{Ingreso máximo.} El ingreso de INN-04 se calcula con el tramo variable máximo mensual (D-14), en el escenario con pago.
\end{itemize}
